\documentclass[12pt]{amsart}

\usepackage{macros}
\usepackage{chngcntr}
\usetikzlibrary{graphs}

\DeclareMathOperator{\Bilin}{Bilin}

\begin{document}

\counterwithout{subsection}{section}

%\section*{Notes on some applications of Young tableaux}
\title{Outline: A paper on Manin triples and spheres}
\author{Surya, Ingmar, Brian}

\maketitle

\section{Introduction}


\section{Local symplectic realizations}

Pavel defines a notion of symplectic realization for a coisotropic morphism. If we have 
\[
C \to X,
\]
then a symplectic realization is a pair of Lagrangian correspondences
\[
\begin{tikzcd}
& & C \ar[dl] \ar[dr] & & \\
& X' \ar[dl] \ar[dr]& & X \ar[dl] \ar[dr] & \\
\star & & Y & & \star,
\end{tikzcd}
\]
so that $Y$ is shifted symplectic,  $X, X' \to Y$ are Lagrangian, and $C \to X' \times_Y X$ is also Lagrangian.

Surya says: A coisotropic morphism should be thought of as a boundary condition for a Poisson degenerate field theory, and an isotropic morphism should be thought of as a boundary condition for a presymplectic degenerate field theory.

The symplectic realization of a Poisson structure is an extension to a Lagrangian morphism that induces that Poisson structure (with appropriate shifts). Thus a coisotropic morphism is a boundary condition for a boundary field theory. To ask for a symplectic realization is to ask for another boundary field theory that makes $C$ into a boundary condition for the sandwich.

This should all be thought of as a coisotropic generalization of the Butson--Yoo universal bulk story.
\section{General package of data: old version}

\subsection{Manin triples}

\begin{dfn} 
A \emph{Manin triple} consists of
\begin{itemize}
\item a Lie algebra $\fd$;
\item an invariant bilinear form on~$\fd$;
\item a pair of transverse Lagrangian subalgebras $\fg^\pm \to \fd$.
\end{itemize}
\end{dfn}

Suppose that $(G^+, G^-, D)$ is a (possibly formal) group triple exponentiating a Manin triple. Then $G^\pm$ inherit Poisson-Lie structures. By a result of Pavel, this data is equivalent to the following. 

\begin{itemize}
\item A 2-shifted symplectic structure on $BD$
\item Lagrangian maps $BG^+, BG^- \to BD$
\item Lagrangian maps $G^\pm \to BG^\mp\times_{BD} BG^\mp$
\end{itemize}

Chern-Simons theories with target $BD$ arise from compactifications of HT theories on $0$-oriented spheres.
\begin{dfn} 
A \emph{$k$-shifted Manin triple} consists of
\begin{itemize}
\item a Lie algebra $\fd$;
\item an invariant bilinear form on~$\fd$ of degree $-k$;
\item a pair of transverse Lagrangian subalgebras $\fg^\pm \to \fd$.
\end{itemize}
\end{dfn}

Given such a datum $B\fd$ is $(2-k)$-shifted symplectic, and $B\fg^\pm\to B\fd$ are Lagrangian. 

Chern-Simons theories with target $B\fd$ arise from compactifications of HT theories on $k$-oriented spheres. Suppose such a $k$-oriented sphere can be written as a fibration of $(k-1)$-oriented spheres over the deRham stack of an interval. Then we may compactify on the total space of the fibration by first compactifying along the fibers and computing a slab compactification in a TNB theory.

Accordingly, note that given a $k$-shifted Manin triple $(\fd , \fg^\pm)$, we get a pair of $k+1$-shifted Manin pairs given by $\Omega B\fg ^\mp \to B \fg^\pm \times _{B \fd} B \fg^\pm$

\subsection{AKSZ-type field theories and corner data from Manin triples}


\subsection{Field theories and universal defects}

\section{Explicit formulas and the second Gelfand--Dickey structure}

\section{Examples}

\section{The model}

Let $X$ be a space and $L$ a subspace. It is a well-known fact in derived geometry that the self-intersection of~$L$ is a shift of its normal bundle:
\[
L \times_X L = N^*_X[-1] L.
\]
(At least, this is true under sufficient assumptions; for example, smooth subvarieties of~$\A^n$.)

We are interested in constructing a model of this computation that will be useful in explicit examples. Down the road, we will want $X$ to be a vector space equipped with a Lie algebra structure. After extending this structure to a model for the derived intersection,  we will study homotopy transfer to an $L_\infty$ algebra structure on~$N^*_X[-1] L = L \oplus (X/L)[-1]$.

Let $I = [0,1]$ denote the unit interval in~$\R$. Consider the following:
\deq{
\Omega[X;L] = \{ f \in \Omega^\bu(I) \otimes_\R X \mid f(t) \in  \Omega^\bu(I) \otimes_\R L \text{ for } t \to 0, ~t \to 1 \} .
}
If one wants, $\Omega[X;L]$ looks like a model for locally constant loops in the relative loop space of $X$ relative to $L$. Note that it is a sub-cdga of $\Omega^\bu(I) \otimes X$.

\begin{prop}
The cohomology of~$\Omega[X;L]$ is $L \oplus (X/L)[-1]$. There is a deformation retract of~$\Omega[X;L]$ onto its cohomology determined by two data:
\begin{itemize}
\item a  positive, monotonically nondecreasing, smooth function $\sigma \in \Omega^0(I)$ with $\sigma = 0$ in a neighborhood of~$0$ and $\sigma=1$ in a neighborhood of~$1$;
\item a right splitting of the sequence
\[
\begin{tikzcd}
0 \ar[r] & L \ar[r, "i"] & X \ar[r, "\pi"] & X/L \ar[r] \ar[l, dashed, bend left, "s"]& 0.
\end{tikzcd}
\]
\end{itemize}
\begin{proof}
We will prove the first claim by constructing the deformation retract explicitly. Let 
\[
\int: \Omega^1(I) \to \Omega^0(I), \qquad \alpha \mapsto \left( t \mapsto \int_0^t \alpha \right)
\] 
denote the indefinite integral starting at $0$, and 
\[
{\int_I} :\Omega^1(I) \to \R, 
\qquad \alpha \mapsto \int_0^1 \alpha
\]
denote the definite integral.
Consider the following diagram:
\[
\begin{tikzcd}[row sep = 4em, column sep = 4em]
L \ar[r, "1\otimes i", shift left] & \Omega[X;L]^0 \ar[d, "d \otimes 1"', shift right] \ar[l, shift left, "\ev_0 \otimes 1"] \\
X/L \ar[r, "d\sigma \otimes s", shift left] & \Omega[X;L]^1 \ar[l, shift left, "\int_I \otimes \pi"] \ar[u, shift right, "h = \int - \,\sigma\cdot s \pi \int_I"']
\end{tikzcd}
\]
Note that the projection in degree $0$ (the map $\ev_0 \otimes 1$) is well-defined due to the boundary conditions imposed on~$\Omega[X;L]$.

Let $f$ denote an element in degree zero. Then $ip(f) = 1 \otimes f(0)$, whereas 
\[
hd(f) = \left (f - f(0) \right) - \sigma \cdot s\pi  \int_I df = f - f(0) - \sigma \cdot s\pi \left( f(1) - f(0) \right).
\]
The last term vanishes due to the boundary conditions on~$f$, so that the required relation holds here.

Let $\alpha$ denote an element in degree one. Then 
\[
ip(\alpha) =d\sigma \otimes s \pi \left( \int_I \alpha \right) ,
\]
whereas
\[
dh(\alpha) = d \int \alpha - d\sigma \cdot s \pi \int_I \alpha = \alpha -  d\sigma \cdot s\pi \int_I \alpha,
\]
so that the condition is satisfied here as well.
\end{proof}
\end{prop}

\section{Manin triples}
Most of this is a poor man's gloss of Pavel's results.

A \emph{Manin triple} is a Lie algebra $\fg$, equipped with the data of an invariant nondegenerate symmetric bilinear pairing and a pair of transverse maximal isotropic subalgebras  denoted $\fg_\pm$.
The pairing gives rise to an identification $\fg_\pm \cong \fg_\mp^\vee$, and we have a collection of maps
\[
\begin{tikzcd}
\fg_+ \ar[r, shift left, "i_+"] & \fg \ar[l, shift left, "\pi_+"] \ar[r, shift left, "\pi_-"] &  \fg_- \ar[l, shift left, "i_-"].
\end{tikzcd}
\]
Note that $\fg_\pm$ are \emph{not} required to be ideals. Thus $i_\pm$ are maps of Lie algebras, whereas $\pi_\pm$ are only maps of vector spaces.
If $\fg_+$ exponentiates to a Lie group $G_+$, then the Manin triple can be used to describe the structure of a Poisson--Lie group on~$G_+$. \textbf{Or is it the other way around? Don't trust this last sentence.}

The pairing equips $B\fg$ with a $2$-shifted symplectic structure, with respect to which $B\fg_\pm$ define two complementary Lagrangians. Thus
\[
B\fg_+ \times_{B\fg} B\fg_+
\]
is a Lagrangian intersection, and is naturally $1$-shifted symplectic. We have seen above that the underlying graded vector space of a minimal $L_\infty$ algebra modeling this intersection takes the form $\fg_+ \oplus \fg_-[-1]$.

Again on general grounds, $\fg_-[-1]$ defines a Lagrangian in the self-intersection, and thereby inherits a $0$-shifted Poisson structure. 
(All of these shifts will change around later.) Our goal is to use the model from above to perform homotopy transfer explicitly, calculate the minimal model of the self-intersection, and match this Poisson structure to the second Adler--Gelfand--Dickey structure. \textbf{I think Pavel's results mean we don't actually have to do this. But oh well.}

$\Omega^\bu(I) \otimes \fg$ inherits a dg Lie algebra structure, and $\Omega[\fg;\fg_+]$ is a sub dg Lie algebra. We apply the previous computation to $\Omega[\fg;\fg_+]$ to see that the cohomology is isomorphic to $\fg_+ \oplus \fg_-[-1]$. Thus we are computing an unshifted Poisson structure on~$B \fg_-[-1]$, which is just the vector space $\fg_-$. To do this, we need to compute the terms in the $L_\infty$ structure on the minimal model that are of the form
\[
\mu_{k+1}(x, y_1, \ldots, y_k),
\]
with $x \in \fg_+$ a degree-zero element and $y_i \in \fg_-$ degree-one elements. The result will be a term in the Poisson bivector of order $(k-1)$ in coefficients. These brackets are given as usual by sums over labeled trees.

Since we are working with the minimal model, there is no differential, and hence no constant-coefficient term. We begin by computing the linear term, which is given by the tree 
\[
\begin{tikzpicture}
\draw (0,0) node[anchor=south] {$x$};
\draw (1,0) node[anchor=south] {$y_1$};
\draw (0,0) -- (1,-2);
\draw (0.5,-1) -- (1,0);
\end{tikzpicture}
\]
The bracket is
\[
[x,y_1] \otimes d\sigma \in \Omega[\fg;\fg_+]^1,
\]
Applying the projection, the output is
\[
\pi_- [x, y_1] \cdot \int_I  d\sigma = \pi_- [x, y_1].
\]
To go further, we apply the homotopy instead, computing the tree
\[
\begin{tikzpicture}
\draw (0,0) node[anchor=south] {$x$};
\draw (1,0) node[anchor=south] {$y_1$};
\draw (2,0) node[anchor=south] {$y_2$};
\draw (0.5, -1.5) node{$h$};
\draw (0,0) -- (1.5,-3);
\draw (0.5,-1) -- (1,0);
\draw (1,-2) -- (2,0);
\end{tikzpicture}
\]
The output is
\[
h( [x,y_1] \otimes d\sigma ) = [x,y_1] \otimes \sigma - \pi_- [x,y_1] \otimes \sigma = \pi_+ [x,y_1] \otimes \sigma.
\]
Going further, we see that the result after bracketing with $i(y_2)$ is
\[
[\pi_+[x,y_1],y_2] \otimes \sigma \, d\sigma.
\]
If we apply the projection, we get 
\[
\frac{1}{2} \pi_- [\pi_+[x,y_1],y_2]
\]
We could also apply the homotopy again, in preparation for computing the next tree:
\[
\begin{tikzpicture}
\draw (0,0) node[anchor=south] {$x$};
\draw (1,0) node[anchor=south] {$y_1$};
\draw (2,0) node[anchor=south] {$y_2$};
\draw (3,0) node[anchor=south] {$y_3$};
\draw (0.5, -1.5) node{$h$};
\draw (1, -2.5) node{$h$};
\draw (0,0) -- (2,-4);
\draw (0.5,-1) -- (1,0);
\draw (1,-2) -- (2,0);
\draw (1.5,-3) -- (3,0);
\end{tikzpicture}
\]
Doing so yields
\[
\frac12 [\pi_+[x,y_1],y_2] \otimes  \sigma^2 - \frac12 \pi_- [\pi_+[x,y_1],y_2] \otimes \sigma.
\]
Bracketing again, we obtain
\[
\frac12 [[\pi_+[x,y_1],y_2], y_3] \otimes \sigma^2\, d\sigma  - \frac12 [\pi_-[\pi_+[x,y_1],y_2], y_3] \otimes \sigma \, d\sigma.
\]
Upon applying the projection, this becomes
\[
\frac16 \pi_- [[\pi_+[x,y_1],y_2], y_3] - \frac14 \pi_- [\pi_-[\pi_+[x,y_1],y_2], y_3].
\]

\end{document}